%=======================================================================
%  CHAPTER 3 — SEARCH TECHNIQUES : SOLVED PROBLEMS
%=======================================================================
\section{Search Techniques \textcolor{sub}{\large -- Solved Problems}}

{\color{sub}\footnotesize Every problem opens with the exam question and the figure
\mk{exactly as the paper prints them}. The OPEN / CLOSED table format is
the one \emph{Insights} \S3.4.2 uses. \gap$\cdot$\gap Every path, cost and minimax value
in this file is recomputed independently by \code{verify.py}.}

\lead{The four things everything here is built on:}
\begin{itemize}
  \item \textbf{1.} $f(n) = g(n) + h(n)$. \textbf{Greedy} orders on $h$ alone,
        \textbf{uniform cost} on $g$ alone, \textbf{A*} on the sum. Every trace below is
        the same loop with a different $f$.
  \item \textbf{2. The loop}: take the best node off OPEN $\rightarrow$ put it on CLOSED
        $\rightarrow$ generate its successors onto OPEN $\rightarrow$ \textbf{re-order
        OPEN} $\rightarrow$ repeat. Stop when the node taken off OPEN is the goal, not
        when the goal is first \emph{generated}.
  \item \textbf{3. If a successor is already known with a larger $g$, replace it} and put
        it back on OPEN. Skipping this returns the wrong answer on both A* graphs in these
        papers, because both heuristics are admissible but \mk{not
        consistent}.
  \item \textbf{4. Alpha-beta}: $\alpha$ only rises, $\beta$ only falls, a child inherits
        both, and you test the cut rule after \emph{every} child. MIN cuts when its value
        $\leq \alpha$; MAX cuts when its value $\geq \beta$.
\end{itemize}

\hr

%=======================================================================
\T{3.N.1 A* on the 2081 Chaitra graph}

\creamq{Explain the A* algorithm on the given figure}{\m{8}}
{\color{sub}\footnotesize \tP{1} \yr{\textbf{81 Ch} \m{8}} \gap$\cdot$\gap
the number \mk{above} each node is $h(n)$, the number on each link is the
step cost}

\begin{asked}{\textbf{81 Ch} \textperiodcentered{} \m{8}}
Explain the A* algorithm on the following figure, where the starting node is S and the
destination is G. (Information in the node represents the cost required to reach the goal
and information in the connecting link represents the cost from one node to the other.)
\figC{c3_81ch_astar.png}{5.3in}
\end{asked}

\sbs{%
\lead{Reading the figure}
\begin{itemize}
  \item $h$: $S{=}17$, $A{=}10$, $B{=}13$, $C{=}4$, $D{=}2$, $E{=}4$, $F{=}1$, $G{=}0$.
  \item Edges: $S$--$A$ 6, $S$--$B$ 5, $S$--$C$ 10, $A$--$E$ 6, $B$--$E$ 6, $B$--$D$ 7,
        $C$--$D$ 6, $E$--$F$ 4, $D$--$F$ 6, $F$--$G$ 3.
  \item There are only four routes from $S$ to $G$, which is worth writing down first as a
        check on the answer:
  \begin{itemize}
    \item $S$-$A$-$E$-$F$-$G$ $= 6{+}6{+}4{+}3 = 19$
    \item $\mathbf{S\text{-}B\text{-}E\text{-}F\text{-}G} = 5{+}6{+}4{+}3 = \mathbf{18}$
    \item $S$-$B$-$D$-$F$-$G$ $= 5{+}7{+}6{+}3 = 21$
    \item $S$-$C$-$D$-$F$-$G$ $= 10{+}6{+}6{+}3 = 25$
  \end{itemize}
  \item \mk{So the answer must come out 18.} If your trace ends at 19 you
        have skipped the reopening step.
\end{itemize}}{%
\lead{Step 1. Start with S on OPEN}
\vspace{-1pt}
\begin{tabularx}{\linewidth}{@{}L{1.1cm}L{0.8cm}L{0.8cm}L{0.8cm}L{1.1cm}X@{}}
\toprule
\hrow \multicolumn{4}{@{}l}{\thead{OPEN}} & \multicolumn{2}{l@{}}{\thead{CLOSED}} \\
\hrow \thead{Node} & \thead{$g$} & \thead{$h$} & \thead{$f$} & \thead{Node} & \thead{Parent} \\
\midrule
S & 0 & 17 & 17 & & \\
\bottomrule
\end{tabularx}
\vspace{3pt}
\lead{Step 2. Expand S}
\vspace{-1pt}
\begin{tabularx}{\linewidth}{@{}L{1.1cm}L{0.8cm}L{0.8cm}L{0.8cm}L{1.1cm}X@{}}
\toprule
\hrow \multicolumn{4}{@{}l}{\thead{OPEN, re-ordered by $f$}} & \multicolumn{2}{l@{}}{\thead{CLOSED}} \\
\hrow \thead{Node} & \thead{$g$} & \thead{$h$} & \thead{$f$} & \thead{Node} & \thead{Parent} \\
\midrule
C & 10 & 4 & \textbf{14} & S & $-$ \\
A & 6 & 10 & 16 & & \\
B & 5 & 13 & 18 & & \\
\bottomrule
\end{tabularx}
\vspace{2pt}
\begin{itemize}
  \item \mk{Note what has already gone wrong for greedy}: $C$ is chosen
        first because $h(C) = 4$ is tiny, yet $C$ lies on the worst route. A* will recover;
        greedy would not.
\end{itemize}}

\penalty0\vspace{2pt}

\sbs{%
\lead{Steps 3--6. Expand C, A, E, F}
\vspace{-1pt}
\begin{tabularx}{\linewidth}{@{}L{1.5cm}L{4.6cm}X@{}}
\toprule
\hrow \thead{Expand} & \thead{Generates} & \thead{OPEN after re-ordering (node/$f$)} \\
\midrule
C \,($f$ 14) & D: $g{=}16$, $f{=}18$ & A/16, B/18, D/18 \\
A \,($f$ 16) & E: $g{=}12$, $f{=}16$ & E/16, B/18, D/18 \\
E \,($f$ 16) & F: $g{=}16$, $f{=}17$ & F/17, B/18, D/18 \\
F \,($f$ 17) & G: $g{=}19$, $f{=}19$ \newline D via F: $g{=}22$, worse, discard & B/18, D/18, G/19 \\
\bottomrule
\end{tabularx}
\vspace{2pt}
\begin{itemize}
  \item $G$ is now \textbf{generated} at cost 19, but it is \mk{not
        expanded}, because $B$ and $D$ both have $f = 18 < 19$. This is exactly
        why the algorithm stops on \emph{removal} from OPEN, not on generation.
\end{itemize}}{%
\lead{\mk{Steps 7--10. Expanding B revises everything}}
\vspace{-1pt}
\begin{tabularx}{\linewidth}{@{}L{1.5cm}L{5.0cm}X@{}}
\toprule
\hrow \thead{Expand} & \thead{Generates} & \thead{OPEN after re-ordering} \\
\midrule
B \,($f$ 18) & E: $g{=}11 < 12$, \textbf{replace}, $f{=}15$ \newline D: $g{=}12 < 16$, \textbf{replace}, $f{=}14$ & D/14, E/15, G/19 \\
D \,($f$ 14) & F via D: $g{=}18 > 15$, discard & E/15, G/19 \\
E \,($f$ 15) & F: $g{=}15 < 16$, \textbf{replace}, $f{=}16$ & F/16, G/19 \\
F \,($f$ 16) & G: $g{=}18 < 19$, \textbf{replace}, $f{=}18$ & G/18 \\
G \,($f$ 18) & \mk{goal removed from OPEN, stop} & \\
\bottomrule
\end{tabularx}
\vspace{2pt}
\begin{itemize}
  \item \textbf{Three replacements}, all through $B$. $B$ looked unpromising ($f = 18$,
        last in the queue) but is on the optimal route.
\end{itemize}}

\penalty0\vspace{2pt}
{\color{acc}\bfseries Answer:} path $\mathbf{S \rightarrow B \rightarrow E \rightarrow F
\rightarrow G}$, total cost $\mathbf{18}$. Nodes expanded in the order
S, C, A, E, F, B, D, E, F, G.

\vspace{2pt}
{\color{sub}\footnotesize \mk{Why this graph is worth understanding rather
than memorising}: $h$ here is \textbf{admissible} (it never exceeds the true
remaining cost) but \textbf{not consistent} $-$ $h(B) = 13$ while
$\text{cost}(B,E) + h(E) = 6 + 4 = 10$. With an inconsistent heuristic a node can be
reached more cheaply after it has been closed, which is why $E$ and $F$ each get expanded
twice. An implementation with a strict closed list and no reopening returns
$S$-$A$-$E$-$F$-$G$ at cost 19, and that is the commonest wrong answer to this question.}

%=======================================================================
\T{3.N.2 A* Vs greedy on the 2080 Chaitra graph}

\creamq{How does A* overcome the problem with greedy best-first search?}{\m{3+4+3}}
{\color{sub}\footnotesize \tP{1} \yr{\textbf{80 Ch} \m{3+4+3}} \gap$\cdot$\gap
\mk{the same graph as \S3.N.1 with different heuristics}, given in a table
rather than on the nodes. Check which you have been given before starting}

\begin{asked}{\textbf{80 Ch} \textperiodcentered{} \m{3+4+3}}
How does A* search overcome the problem associated with Greedy Best First Search? Using the
following figure and table, where the starting node is S and the destination is G, compare
the result of the A* algorithm with greedy search. (The table gives the cost required to
reach the goal; the connecting link gives the cost from one node to the other.)
\figC{c3_80ch_astar.png}{4.3in}
\vspace{2pt}
\begin{center}
\begin{tabular}{|l|c|c|c|c|c|c|c|c|}
\hline
From & S & A & B & C & D & E & F & G \\ \hline
Cost & 15 & 10 & 12 & 5 & 4 & 2 & 1 & 0 \\ \hline
\end{tabular}
\end{center}
\end{asked}

\sbs{%
\lead{Greedy best first: order on $h$ alone}
\vspace{-1pt}
\begin{tabularx}{\linewidth}{@{}L{1.2cm}L{3.4cm}X@{}}
\toprule
\hrow \thead{Expand} & \thead{OPEN after ($h$)} & \thead{CLOSED} \\
\midrule
S & C/5, A/10, B/12 & S \\
C & D/4, A/10, B/12 & S, C \\
D & F/1, A/10, B/12 & S, C, D \\
F & G/0, A/10, B/12 & S, C, D, F \\
G & \mk{goal} & S, C, D, F, G \\
\bottomrule
\end{tabularx}
\vspace{2pt}
\begin{itemize}
  \item \textbf{Path} $S \rightarrow C \rightarrow D \rightarrow F \rightarrow G$,
        \textbf{cost} $10 + 6 + 6 + 3 = \mathbf{25}$.
  \item \mk{That is the worst of the four routes.} Greedy went straight to
        $C$ because $h(C) = 5$ is the smallest number in sight, never noticing that
        \emph{getting} to $C$ costs 10, twice what $B$ costs.
  \item Only 5 nodes expanded, so it is fast. It is fast and wrong.
\end{itemize}}{%
\lead{A*: order on $f = g + h$}
\vspace{-1pt}
\begin{tabularx}{\linewidth}{@{}L{1.2cm}L{1.5cm}X@{}}
\toprule
\hrow \thead{Expand} & \thead{$g$, $h$, $f$} & \thead{OPEN after re-ordering} \\
\midrule
S & 0, 15, 15 & C/15, A/16, B/17 \\
C & 10, 5, 15 & A/16, B/17, D/20 \\
A & 6, 10, 16 & E/14, B/17, D/20 \\
E & 12, 2, 14 & F/17, B/17, D/20 \\
B & 5, 12, 17 & \textbf{E}/13, \textbf{D}/16, F/17, G/19 \\
E & 11, 2, 13 & D/16, F/16, G/19 \\
D & 12, 4, 16 & F/16, G/19 \\
F & 15, 1, 16 & G/18 \\
G & 18, 0, 18 & \mk{goal, stop} \\
\bottomrule
\end{tabularx}
\vspace{2pt}
\begin{itemize}
  \item \textbf{Path} $S \rightarrow B \rightarrow E \rightarrow F \rightarrow G$,
        \textbf{cost} $\mathbf{18}$, the optimum.
\end{itemize}}

\penalty0\vspace{2pt}

\sbs{%
\lead{The comparison, which is what the 3 marks are for}
\vspace{-1pt}
\begin{tabularx}{\linewidth}{@{}L{2.0cm}L{3.1cm}X@{}}
\toprule
\hrow & \thead{Greedy best first} & \thead{A*} \\
\midrule
$f(n)$ & $h(n)$ & $g(n) + h(n)$ \\
Path found & $S$-$C$-$D$-$F$-$G$ & $S$-$B$-$E$-$F$-$G$ \\
Cost & \mk{25}, the worst route & \mk{18}, the optimum \\
Nodes expanded & 5 & 9 \\
Optimal & \xmark no & \checkmark yes, $h$ is admissible \\
Complete & \xmark no & \checkmark yes \\
\bottomrule
\end{tabularx}
\vspace{2pt}
\begin{itemize}
  \item \mk{The trade in one line}: A* expanded almost twice as many nodes
        and returned a path 28\,\% cheaper. Greedy buys speed with correctness.
\end{itemize}}{%
\lead{Where greedy went wrong, precisely}
\begin{itemize}
  \item At the very first step both algorithms rank $C$ best, but for different reasons.
        Greedy sees $h(C) = 5$ and \textbf{commits}; A* sees $f(C) = 10 + 5 = 15$, ties
        with nothing better, expands it, and then \mk{keeps $B$ on the
        table}.
  \item Once greedy has moved to $C$ it never reconsiders. It has \textbf{no memory of
        cost paid}, so nothing can ever make $C$ look like a mistake.
  \item \textbf{The general statement to close with}: a heuristic estimates the
        \emph{remaining} journey. Judging a node on that alone is judging a route by how
        it ends, ignoring how much it cost to get there. Adding $g(n)$ restores the whole
        journey to the comparison, and that is the entire difference between the two
        algorithms.
\end{itemize}}

%=======================================================================
\T{3.N.3 Best first search on the 2076 Bhadra tree}

\creamq{Show the open list, closed list and path at each step}{\m{5}}
{\color{sub}\footnotesize \yr{\textbf{\textit{76 Bh}} \m{5}} \gap$\cdot$\gap
1 paper, 5 marks, so no tier chip \gap$\cdot$\gap
the question names the \mk{three things it wants tabulated}, so tabulate
exactly those three}

\begin{asked}{\textbf{\textit{76 Bh}} \textperiodcentered{} \m{5}}
Consider a search tree for a path-finding problem as shown. Assuming node A to be the root
node and node I to be the goal node, explain how Best First Search reaches the goal by
showing changes in the open list, closed list, and path at each step of the search.
\figC{c3_76bh_bfs.png}{5.0in}
\end{asked}

\sbs{%
\lead{Reading the figure}
\begin{itemize}
  \item Tree: $A \rightarrow B, C, D$; \quad $B \rightarrow E, F$; \quad
        $D \rightarrow G, H$; \quad $G \rightarrow I, J$. \quad $I$ (the square) is the
        goal.
  \item Values on the edges, taken as the estimate for the child node:
        $B{=}3$, $C{=}6$, $D{=}1$, $E{=}6$, $F{=}5$, $G{=}4$, $H{=}6$, $I{=}1$, $J{=}2$.
  \item \textbf{Best first} expands the OPEN node with the smallest value, so the OPEN
        list is re-ordered after every expansion.
\end{itemize}
\lead{The trace}
\vspace{-1pt}
\begin{tabularx}{\linewidth}{@{}L{0.75cm}L{1.1cm}XL{2.4cm}@{}}
\toprule
\hrow \thead{Step} & \thead{Expand} & \thead{OPEN after (node/value)} & \thead{CLOSED} \\
\midrule
1 & A & D/1, B/3, C/6 & A \\
2 & D & B/3, G/4, C/6, H/6 & A, D \\
3 & B & G/4, F/5, C/6, E/6, H/6 & A, D, B \\
4 & G & \textbf{I}/1, J/2, F/5, C/6, E/6, H/6 & A, D, B, G \\
5 & I & \mk{goal reached} & A, D, B, G, I \\
\bottomrule
\end{tabularx}}{%
\lead{Path at each step}
\vspace{-1pt}
\begin{tabularx}{\linewidth}{@{}L{0.75cm}X@{}}
\toprule
\hrow \thead{Step} & \thead{Path to the node just expanded} \\
\midrule
1 & A \\
2 & A $\rightarrow$ D \\
3 & A $\rightarrow$ B \\
4 & A $\rightarrow$ D $\rightarrow$ G \\
5 & \textbf{A $\rightarrow$ D $\rightarrow$ G $\rightarrow$ I} \\
\bottomrule
\end{tabularx}
\vspace{3pt}
{\color{acc}\bfseries Answer:} path $\mathbf{A \rightarrow D \rightarrow G \rightarrow I}$,
cost $1 + 4 + 1 = \mathbf{6}$.
\vspace{2pt}
\begin{itemize}
  \item \mk{Step 3 is the one to explain.} After expanding $D$, the best
        node on OPEN is $B$ at 3, not $G$ at 4 $-$ so best first
        \textbf{jumps back across the tree} to a different branch before returning to
        $G$. That switching is precisely what distinguishes best first from depth first,
        and it is where the marks are.
  \item $C$, $E$, $F$, $H$ and $J$ are \textbf{never expanded}. Five of the ten nodes are
        never touched, which is the point of using a heuristic at all.
  \item If the numbers are read as \emph{edge costs} and accumulated instead, the same
        path comes out at the same cost of 6. Either reading is defensible here; say which
        you are using.
\end{itemize}}

%=======================================================================
\T{3.N.4 Minimax and alpha-beta on the 2079 Chaitra tree}

\creamq{Apply minimax with alpha-beta pruning to the given tree}{\m{2+6}}
{\color{sub}\footnotesize \tP{1} \yr{\textbf{79 Ch} \m{2+6}} \gap$\cdot$\gap
squares are MAX, circles are MIN, and \mk{three of the MIN nodes have only
one child}, which is easy to misread}

\begin{asked}{\textbf{79 Ch} \textperiodcentered{} \m{2+6}}
What is minmax search? Apply minmax search along with alpha-beta pruning for the following
tree (max node at top).
\figC{c3_79ch_minmax.png}{5.0in}
\end{asked}

\sbs{%
\lead{Reading the tree}
\begin{itemize}
  \item Five levels: \textbf{MAX} root, \textbf{MIN}, \textbf{MAX}, \textbf{MIN}, then the
        12 leaves.
  \item Left MIN node $m_1$ over MAX nodes $s_1, s_2$; right MIN node $m_2$ over
        $s_3, s_4$.
  \item The MIN nodes above the leaves, left to right, take:
        $\{3,9\}$, $\{2,7\}$, $\{2\}$, $\{5,0\}$, $\{2,5\}$, $\{8\}$, $\{3,14\}$.
        \mk{Three of them have a single child}: count the lines in the
        figure, not the leaves.
\end{itemize}
\lead{Step 1: MIN over the leaves}
\begin{itemize}
  \item $c_1 = \min(3,9) = 3$ \quad $c_2 = \min(2,7) = 2$ \quad $c_3 = 2$
  \item $c_4 = \min(5,0) = 0$ \quad $c_5 = \min(2,5) = 2$ \quad $c_6 = 8$
  \item $c_7 = \min(3,14) = 3$
\end{itemize}
\lead{Step 2: MAX}
\begin{itemize}
  \item $s_1 = \max(3,2) = 3$ \quad $s_2 = \max(2,0) = 2$
  \item $s_3 = \max(2,8) = 8$ \quad $s_4 = 3$
\end{itemize}
\lead{Step 3: MIN, then the root}
\begin{itemize}
  \item $m_1 = \min(3,2) = 2$ \quad $m_2 = \min(8,3) = 3$
  \item \textbf{root} $= \max(2,3) = \mathbf{3}$, so MAX plays the \textbf{right} branch.
\end{itemize}}{%
\lead{Alpha-beta, left to right}
\vspace{-1pt}
\begin{tabularx}{\linewidth}{@{}L{1.1cm}L{1.9cm}X@{}}
\toprule
\hrow \thead{Node} & \thead{$(\alpha, \beta)$ on entry} & \thead{What happens} \\
\midrule
$c_1$ & $(-\infty, \infty)$ & leaves 3, 9 $\rightarrow$ 3 \\
$s_1$ & $\alpha \leftarrow 3$ & \\
$c_2$ & $(3, \infty)$ & leaf 2 gives value 2; $2 \leq \alpha = 3$ $\rightarrow$ \mk{PRUNE leaf 7} \\
$m_1$ & $\beta \leftarrow 3$ & $s_1 = 3$ \\
$c_3$ & $(-\infty, 3)$ & leaf 2 $\rightarrow$ 2 \\
$c_4$ & $(2, 3)$ & leaves 5, 0 $\rightarrow$ 0; no cut, 5 $>\alpha$ when tested \\
$m_1$ & & $= \min(3, 2) = 2$; root $\alpha \leftarrow 2$ \\
$c_5$ & $(2, \infty)$ & leaf 2 gives 2; $2 \leq \alpha = 2$ $\rightarrow$ \mk{PRUNE leaf 5} \\
$c_6$ & $(2, \infty)$ & leaf 8 $\rightarrow$ 8, so $s_3 = 8$ \\
$m_2$ & $\beta \leftarrow 8$ & \\
$c_7$ & $(2, 8)$ & leaves 3, 14 $\rightarrow$ 3, so $s_4 = 3$ \\
root & & $m_2 = \min(8,3) = 3$; $\max(2,3) = \mathbf{3}$ \\
\bottomrule
\end{tabularx}}

\penalty0\vspace{2pt}
{\color{acc}\bfseries Answer:} minimax value $\mathbf{3}$; MAX plays the right branch.
\textbf{2 of the 12 leaves are pruned}: the \textbf{7} under $c_2$, and the \textbf{5}
under $c_5$. Ten leaves are examined.

\vspace{2pt}
{\color{sub}\footnotesize \mk{Note the equality in the second cut.} At
$c_5$ the value is 2 and $\alpha$ is also 2. A MIN node cuts on $\textbf{value} \leq
\alpha$, not $<$, because MAX already has a way to score 2 elsewhere and will never
\emph{prefer} this branch. Using a strict $<$ here is the commonest slip and it costs you
that pruned leaf.}

%=======================================================================
\T{3.N.5 Alpha-beta on the 2076 Baishakh tree}

\creamq{Discuss alpha-beta pruning and find the minimax value}{\m{4+4}}
{\color{sub}\footnotesize \tP{1} \yr{76 Ba \m{4+4}} \gap$\cdot$\gap
a small tree that prunes \mk{a much larger fraction} than \S3.N.4, and the
reason why is the lesson}

\begin{asked}{76 Ba \textperiodcentered{} \m{4+4}}
Discuss the alpha-beta pruning algorithm. Find the min-max value using this concept in the
following tree.
\figC{c3_76ba_minmax.png}{5.0in}
\end{asked}

\sbs{%
\lead{The trace}
\begin{itemize}
  \item MAX root over three MIN nodes, each with three leaves:
        $\{3,5,10\}$, $\{2,8,19\}$, $\{2,7,3\}$.
\end{itemize}
\vspace{-1pt}
\begin{tabularx}{\linewidth}{@{}L{1.3cm}L{1.7cm}X@{}}
\toprule
\hrow \thead{Node} & \thead{$(\alpha, \beta)$} & \thead{What happens} \\
\midrule
$m_1$ & $(-\infty, \infty)$ & 3 $\rightarrow$ 5 $\rightarrow$ 10, value $= 3$.
  \mk{Nothing can be pruned here}, because $\alpha$ is still $-\infty$ \\
root & $\alpha \leftarrow 3$ & \\
$m_2$ & $(3, \infty)$ & leaf 2 gives 2; $2 \leq 3$ $\rightarrow$ \mk{PRUNE 8 and 19} \\
$m_3$ & $(3, \infty)$ & leaf 2 gives 2; $2 \leq 3$ $\rightarrow$ \mk{PRUNE 7 and 3} \\
root & & $\max(3, 2, 2) = \mathbf{3}$ \\
\bottomrule
\end{tabularx}
\vspace{2pt}
\begin{itemize}
  \item \textbf{Minimax value 3}, MAX plays the \textbf{leftmost} branch.
  \item \textbf{4 of the 9 leaves pruned} (8, 19, 7, 3); five examined:
        3, 5, 10, 2, 2.
\end{itemize}}{%
\lead{\mk{Why this tree prunes so much better}}
\begin{itemize}
  \item \S3.N.4 pruned 2 leaves in 12, about 17\,\%. This one prunes 4 in 9, about 44\,\%.
        The algorithm is identical; the \textbf{move ordering} is not.
  \item Here the \textbf{best branch happens to be examined first}. That sets $\alpha = 3$
        immediately, and every later MIN node that finds any child below 3 can stop at
        once. Both of them do, on their very first leaf.
  \item \textbf{The general rule}: alpha-beta is only as good as the order in which moves
        are examined.
  \begin{itemize}
    \item \textbf{Best case} (best move first everywhere): $O(b^{m/2})$, so the same time
          searches \textbf{twice as deep}.
    \item \textbf{Worst case} (best move last everywhere): $O(b^m)$, i.e.\
          \mk{no saving at all} $-$ alpha-beta never prunes more than minimax
          examines, and can prune nothing.
  \end{itemize}
  \item This is why real game programs spend effort on \textbf{move ordering heuristics}
        before they spend it on searching deeper.
\end{itemize}
\lead{Two things to say when asked to ``discuss'' it}
\begin{itemize}
  \item \mk{It returns exactly the minimax value.} Pruning removes work,
        never accuracy. Say this first.
  \item A pruned branch is one whose value \textbf{cannot change the root}, because a
        player higher up already has a better guaranteed alternative.
\end{itemize}}
